\documentclass[12pt,a4paper]{article}
\usepackage[utf8]{inputenc}
\usepackage[T2A]{fontenc}
\usepackage[russian]{babel}
\usepackage{textcomp}
\usepackage{amsmath, amssymb, amsthm}
\usepackage{graphicx}
\usepackage{hyperref}
\usepackage{listings}
\usepackage{xcolor}
\usepackage{geometry}
\usepackage{float}
\usepackage{caption}
\usepackage{subcaption}
\usepackage{enumitem}
\usepackage{titlesec}
\usepackage{fancyhdr}
\usepackage{array}
\usepackage{booktabs}

% Настройка полей
\geometry{
    a4paper,
    left=2.5cm,
    right=1.5cm,
    top=2cm,
    bottom=2cm
}

% Настройка заголовков
\titleformat{\section}{\Large\bfseries}{\thesection}{1em}{}
\titleformat{\subsection}{\large\bfseries}{\thesubsection}{1em}{}
\titleformat{\subsubsection}{\normalsize\bfseries}{\thesubsubsection}{1em}{}

% Настройка кода
\lstset{
    language=Python,
    basicstyle=\ttfamily\small,
    keywordstyle=\color{blue},
    commentstyle=\color{green!40!black},
    stringstyle=\color{red},
    numbers=left,
    numberstyle=\tiny,
    breaklines=true,
    frame=single,
    captionpos=b,
    showstringspaces=false,
    tabsize=4
}

% Настройка колонтитулов
\setlength{\headheight}{15pt}
\pagestyle{fancy}
\fancyhf{}
\fancyhead[L]{Лабораторная работа №2: Развитие AI-сервиса}
\fancyhead[R]{\thepage}
\fancyfoot[C]{}

\begin{document}

\begin{titlepage}
    \begin{center}
        \textbf{МИНИСТЕРСТВО НАУКИ И ВЫСШЕГО ОБРАЗОВАНИЯ РОССИЙСКОЙ ФЕДЕРАЦИИ}
        \vspace{0.5cm}

        \textbf{Федерального государственного автономного образовательного учреждения \\ высшего образования}
        \vspace{0.3cm}

        \textbf{«НАЦИОНАЛЬНЫЙ ИССЛЕДОВАТЕЛЬСКИЙ ТЕХНОЛОГИЧЕСКИЙ УНИВЕРСИТЕТ МИСИС»}
        \vspace{2cm}

        \textbf{Институт компьютерных наук}
        \vspace{0.3cm}

        \textbf{Кафедра Инженерной кибернетики}
        \vspace{1cm}

        \textbf{Лабораторная работа №2}
        \vspace{0.3cm}

        \textbf{по дисциплине «Современные интеллектуальные сетевые сервисы»}
        \vspace{1.5cm}

        \huge{\textbf{Развитие AI-сервиса: \\ сквозная идентификация запросов, \\ асинхронная обработка и интеграция LLM}}
        \vspace{2cm}
    \end{center}

    \begin{flushright}
        \textbf{Выполнил:} \\
        Студент группы МПИ 25-1-2\\
        Д.А. Комарова \\
        \vspace{0.5cm}
        \textbf{Проверил:} \\
        Преподаватель \\
        М.Е. Корнет \\
        \vspace{2cm}
        \textbf{Оценка:} \hrulefill
    \end{flushright}

    \begin{center}
        \vfill
        Москва, 2026
    \end{center}
\end{titlepage}

\tableofcontents
\newpage

\section{Цель работы}

Цель практической работы — развить ранее созданный REST API сервиса агроскоринга и приблизить его архитектуру к реальному промышленному AI-сервису. В ходе работы необходимо изучить и реализовать механизмы, которые возникают при эксплуатации сервиса: сквозную идентификацию запросов с помощью \texttt{request\_id}, использование промежуточных слоёв (middleware), управление очередью запросов и пакетным инференсом, хранение результатов в базе данных SQLite, асинхронное взаимодействие с внешними сервисами и подключение большой языковой модели через HTTP API.

\section{Задачи работы}

В ходе лабораторной работы необходимо реализовать:
\begin{enumerate}
    \item Модификацию API для назначения уникального \texttt{request\_id}.
    \item Middleware для извлечения/генерации \texttt{request\_id} и измерения времени выполнения.
    \item Два режима инференса: одиночный и пакетный.
    \item Очередь запросов на базе \texttt{asyncio.Queue}.
    \item Динамическое формирование пакета (dynamic batching).
    \item Эксплуатационные ограничения (размер batch, время ожидания, размер очереди, таймаут).
    \item Механизм backpressure (возврат HTTP 429 при переполнении).
    \item Подключение SQLite для постоянного хранения.
    \item Асинхронную работу с БД через \texttt{aiosqlite}.
    \item Запись, чтение, поиск по ID и пагинацию результатов.
    \item Подключение внешней LLM через HTTP API.
    \item Асинхронный HTTP-клиент \texttt{httpx.AsyncClient}.
    \item Endpoint для запроса к LLM.
    \item Обработку ошибок внешних зависимостей (400, 404, 409, 422, 429, 500, 502, 503, 504).
    \item Тестирование и анализ полного пути запроса.
\end{enumerate}

\section{Краткая теоретическая часть}

\textbf{Request\_id} используется для сквозной идентификации запроса в микросервисной архитектуре. Middleware перехватывает запрос, проверяет наличие заголовка \texttt{X-Request-ID}, и если его нет — генерирует UUID, сохраняя в \texttt{request.state}.

\textbf{Batching} повышает throughput за счёт группировки запросов. \textit{Explicit batching} — клиент сам присылает пакет. \textit{Dynamic batching} — сервер накапливает одиночные запросы в очереди и отправляет их в модель пачкой при достижении лимита \texttt{MAX\_BATCH\_SIZE} или по истечении \texttt{MAX\_BATCH\_WAIT\_MS}.

\textbf{Backpressure} ограничивает размер очереди. Если она переполнена, сервис возвращает 429 Too Many Requests, защищая память от перегрузки и деградации latency.

\textbf{SQLite + aiosqlite} позволяют асинхронно работать с файловой БД без блокировки event loop. Режим WAL (Write-Ahead Logging) позволяет читателям и писателям не блокировать друг друга.

\textbf{LLM интеграция} осуществляется через \texttt{httpx.AsyncClient} с таймаутом. Если внешний сервис недоступен, FastAPI возвращает 502/503/504.

\section{Результаты выполнения заданий}

\subsection{Задание 1-2. Сквозной request\_id и Middleware}
Реализовано middleware \texttt{request\_context\_middleware}, которое перехватывает каждый запрос. Оно извлекает \texttt{X-Request-ID} из заголовков или генерирует новый UUID с помощью \texttt{uuid.uuid4()}, сохраняет его в \texttt{request.state.request\_id}, замеряет время выполнения через \texttt{time.perf\_counter()} и добавляет заголовки \texttt{X-Request-ID} и \texttt{X-Process-Time} в ответ.

\subsection{Задание 3. Обработка ошибок}
Настроены обработчики для возврата кодов 400 (неизвестный регион), 404 (прогноз не найден), 409 (несколько прогнозов для одного request\_id), 422 (ошибки Pydantic при неверных типах или размере batch), 429 (переполнение очереди), 502/503/504 (ошибки внешней LLM и таймауты).

\subsection{Задание 4-6. SQLite и чтение результатов}
Создан модуль \texttt{database.py} с использованием \texttt{aiosqlite}. Таблица \texttt{predictions} содержит все необходимые поля. \texttt{POST /predict} сохраняет результат в БД. \texttt{GET /predictions} и \texttt{GET /predictions/\{request\_id\}} читают данные из SQLite. Если \texttt{request\_id} относится к batch и содержит несколько записей, возвращается 409.

\textbf{Ответ на вопрос:} Блокирующий \texttt{sqlite3.execute()} останавливает весь event loop на время работы с диском, что приводит к таймаутам у других клиентов. \texttt{aiosqlite} выносит операции в отдельный поток.

\subsection{Задание 7. Пагинация}
Реализован endpoint \texttt{GET /predictions/by-weekday}, который с помощью \texttt{LIMIT} и \texttt{OFFSET} возвращает порцию записей для выбранного дня недели, а также метаинформацию (текущая страница, всего страниц, всего элементов).

\subsection{Задание 8. Explicit batching}
Реализован \texttt{POST /batch-predict}, принимающий список объектов. Размер списка ограничен 8 элементами. Обработка происходит в цикле, результаты сохраняются в БД.

\textbf{Ответ на вопрос:} Один \texttt{executemany()} создаёт одну транзакцию и один вызов fsync на диск. Восемь отдельных \texttt{INSERT + COMMIT} вызовут восемь синхронизаций диска, что в сотни раз медленнее.

\subsection{Задание 9-11. Очередь, backpressure и timeout}
Создан \texttt{queue\_manager.py} с \texttt{asyncio.Queue}. Воркер накапливает запросы до 8 штук или ждёт 50 мс, затем обрабатывает их. Реализован endpoint \texttt{POST /predict-queued}. Если очередь заполнена (100 элементов), возвращается 429. Время ожидания ответа от модели ограничено \texttt{INFERENCE\_TIMEOUT\_SECONDS} через \texttt{asyncio.wait\_for}. При превышении возвращается 504.

\textbf{Ответ на вопрос:} При неограниченной очереди память будет расти бесконечно, пока сервер не упадёт с OOM. Latency также будет расти бесконечно, так как запросы будут ждать часами.

\subsection{Задание 12-14. Внешняя LLM}
Создан \texttt{llm\_service.py}, использующий \texttt{httpx.AsyncClient} для запросов к Vireonix. Реализован \texttt{POST /llm/chat} для прямого запроса. Реализован \texttt{POST /predict-with-llm}, который сначала считает риск, а затем просит LLM объяснить результат. Если LLM недоступна, возвращается только результат скоринга с пометкой об ошибке.

\section{Скриншоты работы приложения}

\begin{figure}[H]
    \centering
    \includegraphics[width=\textwidth]{01_swagger.png}
    \caption{Общий вид Swagger UI}
\end{figure}

\begin{figure}[H]
    \centering
    \includegraphics[width=\textwidth]{02_predict.png}
    \caption{Выполнение POST /predict}
\end{figure}

\begin{figure}[H]
    \centering
    \includegraphics[width=\textwidth]{03_db_read.png}
    \caption{Чтение результатов из SQLite (GET /predictions)}
\end{figure}

\begin{figure}[H]
    \centering
    \includegraphics[width=\textwidth]{04_batch.png}
    \caption{Явный батчинг (POST /batch-predict)}
\end{figure}

\begin{figure}[H]
    \centering
    \includegraphics[width=\textwidth]{05_queue_status.png}
    \caption{Статус очереди (GET /queue-status)}
\end{figure}

\begin{figure}[H]
    \centering
    \includegraphics[width=\textwidth]{06_pagination.png}
    \caption{Пагинация по дню недели}
\end{figure}

\begin{figure}[H]
    \centering
    \includegraphics[width=\textwidth]{07_llm_chat.png}
    \caption{Запрос к внешней LLM (504 Timeout)}
\end{figure}

\begin{figure}[H]
    \centering
    \includegraphics[width=\textwidth]{08_predict_llm.png}
    \caption{Инференс с объяснением от LLM}
\end{figure}

\section{Таблица тест-кейсов}

\begin{table}[H]
\centering
\caption{Таблица тест-кейсов с ожидаемыми HTTP-кодами}
\label{tab:test_cases}
\begin{tabular}{|c|p{5cm}|c|c|c|}
\hline
\textbf{№} & \textbf{Запрос} & \textbf{Ожидаемый} & \textbf{Полученный} & \textbf{Результат} \\
\hline
1  & GET /health & 200 & 200 & Успех \\
2  & корректный POST /predict & 200 & 200 & Успех \\
3  & area\_ha = -100 & 422 & 422 & Успех \\
4  & неизвестный регион & 400 & 400 & Успех \\
5  & существующий request\_id & 200 & 200 & Успех \\
6  & неизвестный request\_id & 404 & 404 & Успех \\
7  & корректный POST /batch-predict & 200 & 200 & Успех \\
8  & batch больше MAX\_BATCH\_SIZE & 422 & 422 & Успех \\
9  & POST /predict-queued & 200 & 200 & Успех \\
10 & переполненная очередь & 429 & 429 & Успех \\
11 & GET /predictions/by-weekday & 200 & 200 & Успех \\
12 & page = -1 & 422 & 422 & Успех \\
13 & неизвестный weekday & 400 & 400 & Успех \\
14 & корректный POST /llm/chat & 200 & 504 & Ошибка LLM \\
15 & недоступный LLM & 502/503/504 & 504 & Ожидаемо \\
16 & корректный POST /predict-with-llm & 200 & 200 & Успех \\
\hline
\end{tabular}
\end{table}

\section{Таблица API-контракта}

\begin{table}[H]
\centering
\caption{Таблица API-контракта}
\label{tab:api_contract}
\begin{tabular}{|p{3.2cm}|p{1.4cm}|p{2.4cm}|p{3cm}|p{1.2cm}|p{2.4cm}|}
\hline
\textbf{Endpoint} & \textbf{Method} & \textbf{Input} & \textbf{Output} & \textbf{Success} & \textbf{Возможные ошибки} \\
\hline
/health & GET & --- & status & 200 & 500 \\
/predict & POST & FarmRequest & PredictionResponse & 200 & 400, 422, 503, 504 \\
/batch-predict & POST & BatchRequest & BatchPredictionResponse & 200 & 400, 422, 503, 504 \\
/predict-queued & POST & FarmRequest & PredictionResponse & 200 & 400, 422, 429, 503, 504 \\
/predictions & GET & query & list & 200 & 400, 422, 503 \\
/predictions/by-weekday & GET & weekday, page, per\_page & paginated response & 200 & 400, 404, 422 \\
/llm/chat & POST & LLMChatRequest & LLMChatResponse & 200 & 422, 502, 503, 504 \\
/predict-with-llm & POST & FarmRequest & PredictWithLLMResponse & 200 & 400, 422, 502, 503, 504 \\
\hline
\end{tabular}
\end{table}

\section{Ссылка на исходный код}

Исходный код приложения доступен в репозитории: \url{https://github.com/dashakomarova2002-commits/HW_2_SISS}

\section{Вывод}

В ходе работы мы расширили архитектуру AI-сервиса агроскоринга. \texttt{request\_id} обеспечивает сквозную трассировку, позволяя связать логи, ошибки и записи в БД с конкретным клиентским запросом. Middleware изолирует эту логику от бизнес-кода. Batching (явный и динамический) позволяет повысить пропускную способность за счёт группировки запросов к модели. Backpressure (429) и таймауты (504) защищают сервис от перегрузки и зависаний. База данных SQLite (через \texttt{aiosqlite}) обеспечивает надежное постоянное хранение результатов без блокировки event loop. Интеграция с LLM через \texttt{httpx.AsyncClient} позволяет обогащать результаты текстовыми объяснениями, при этом сервис устойчив к сбоям внешнего провайдера, корректно возвращая коды 502/503/504.

\section*{Приложение А: Исходный код}

\begin{lstlisting}[caption={main.py}]
import uuid
import time
import asyncio
import logging
from contextlib import asynccontextmanager
from fastapi import FastAPI, Request, Response, HTTPException, Query, status

from schemas import (
    FarmRequest, PredictionResponse, BatchRequest, BatchPredictionResponse,
    PaginatedResponse, QueueStatusResponse, LLMChatRequest, LLMChatResponse,
    PredictWithLLMResponse
)
from services import build_prediction, is_valid_region
from database import (
    init_db, save_one_prediction, get_predictions_by_req_id,
    get_predictions_list, get_predictions_by_weekday_paginated
)
from queue_manager import (
    queue, worker, QUEUE_MAX_SIZE, MAX_BATCH_SIZE, MAX_BATCH_WAIT_MS,
    INFERENCE_TIMEOUT_SECONDS
)
from llm_service import call_llm

logging.basicConfig(level=logging.INFO, format="%(asctime)s [%(levelname)s] %(message)s")
logger = logging.getLogger(__name__)

@asynccontextmanager
async def lifespan(app: FastAPI):
    await init_db()
    task = asyncio.create_task(worker())
    yield
    task.cancel()

app = FastAPI(
    title="Agro Scoring API",
    description="API for agricultural risk scoring (Lab 2)",
    version="2.0.0",
    lifespan=lifespan
)

@app.middleware("http")
async def request_context_middleware(request: Request, call_next):
    request_id = request.headers.get("X-Request-ID") or str(uuid.uuid4())
    request.state.request_id = request_id
    start_time = time.perf_counter()
    logger.info(f"[{request_id}] START {request.method} {request.url.path}")

    try:
        response: Response = await call_next(request)
    except Exception as e:
        logger.error(f"[{request_id}] ERROR {str(e)}")
        response = Response(
            content=f'{{"detail": "Internal server error", "request_id": "{request_id}"}}',
            status_code=500,
            media_type="application/json"
        )

    process_time = time.perf_counter() - start_time
    response.headers["X-Request-ID"] = request_id
    response.headers["X-Process-Time"] = f"{process_time:.4f}"
    logger.info(f"[{request_id}] END {request.method} {request.url.path} - {response.status_code} in {process_time:.4f}s")
    return response

@app.get("/health")
async def health(request: Request):
    return {"status": "ok", "request_id": request.state.request_id}

@app.post("/predict", response_model=PredictionResponse, status_code=status.HTTP_200_OK)
async def predict(request_body: FarmRequest, request: Request):
    request_id = request.state.request_id
    if not is_valid_region(request_body.region):
        raise HTTPException(status_code=400, detail="Unknown region")

    result = build_prediction(request_body, request_id)
    await save_one_prediction(result, source="single")
    return result

# ... (other endpoints are omitted for brevity)
\end{lstlisting}

\end{document}